\ifdefined\gkver\else\def\gkver{student}\fi
\documentclass[\gkver]{gaokaozhenti}
\newcommand{\ccnum}[1]{{\fontspec{Noto Sans CJK JP}#1}}
\begin{document}

\chapter{2002年全国理综旧课程}
%% paperType: 理综物理部分

\begin{enumerate}
\item
%% number: 15
%% paperName: 2002年全国理综旧课程第 15 题 6 分
%% typeId: 1
%% type: 单选题
%% chapter: 原子和原子核
%% point: 原子核的组成
%% method: 无
%% score: 6
%% degree: 650
%% duplicateId: 0
%% body:
目前普通认为，质子和中子都是由被称为 u 夸克和 d 夸克的两类夸克组成。u 夸克带电量为 $\frac{2}{3}e$，d 夸克带电量为 $-\frac{1}{3}e$，$e$ 为基元电荷。下列论断可能正确的是\xzanswer{B}

\fourchoices[answer=B]
{质子由 1 个 u 夸克和 1 个 d 夸克组成，中子由 1 个 u 夸克和 2 个 d 夸克组成}
{质子由 2 个 u 夸克和 1 个 d 夸克组成，中子由 1 个 u 夸克和 2 个 d 夸克组成}
{质子由 1 个 u 夸克和 2 个 d 夸克组成，中子由 2 个 u 夸克和 1 个 d 夸克组成}
{质子由 2 个 u 夸克和 1 个 d 夸克组成，中子由 1 个 u 夸克和 1 个 d 夸克组成}

%% answer: B
%% memo:
\memoanswer{
由质子带一个单位正电荷，中子不带电，设质子中 u 夸克、d 夸克个数分别是 $x$、$y$（$x$、$y$ 取正整数），则 $x\times\frac{2}{3}+y\times\left(-\frac{1}{3}\right)=1$，解得 $x=2$、$y=1$；设中子中 u 夸克、d 夸克个数分别是 $m$、$n$（$m$、$n$ 取正整数），$m\times\left(\frac{2}{3}\right)+n\times\left(-\frac{1}{3}\right)=0$，解得 $m=1$、$n=2$，故 B 正确。
}

\item
%% number: 16
%% paperName: 2002年全国理综旧课程第 16 题 6 分
%% typeId: 1
%% type: 单选题
%% chapter: 运动和力
%% point: 动量和能量
%% method: 无
%% score: 6
%% degree: 550
%% duplicateId: 0
%% body:
在光滑水平地面上有两个相同的弹性小球 A、B，质量都为 $m$。现 B 球静止，A 球向 B 球运动，发生正碰。已知碰撞过程中总机械能守恒，两球压缩最紧时的弹性势能为 $E_{p}$，则碰前 A 球的速度等于\xzanswer{C}

\fourchoices[answer=C]
{$\sqrt{\frac{E_{p}}{m}}$}
{$\sqrt{\frac{2E_{p}}{m}}$}
{$2\sqrt{\frac{E_{p}}{m}}$}
{$2\sqrt{\frac{2E_{p}}{m}}$}

%% answer: C
%% memo:
\memoanswer{
设碰撞前 A 球的速度为 $v_{0}$，两个弹性小球发生正碰，当两者速度相同时，弹性势能最大，由动量守恒得 $mv_{0}=2mv$，由机械能守恒得 $E_{p}=\frac{1}{2}mv_{0}^{2}-\frac{1}{2}\cdot 2mv^{2}$，解得 $v_{0}=2\sqrt{\frac{E_{p}}{m}}$，故 C 正确。
}

\item
%% number: 17
%% paperName: 2002年全国理综旧课程第 17 题 6 分
%% typeId: 1
%% type: 单选题
%% chapter: 电路
%% point: 电容
%% method: 无
%% score: 6
%% degree: 650
%% duplicateId: 0
%% body:
图中 EF、GH 为平行的金属导轨，其电阻可不计，R 为电阻器，G 为电容器，AB 为可在 EF 和 GH 上滑动的导体横杆。有均匀磁场垂直与导轨平面。若用 $I_{1}$ 和 $I_{2}$ 分别表示图中该处导线中的电流，则当横杆 AB\xzanswer{D}

\onepicture[w=6.0cm]{figs/17.png}

\fourchoices[answer=D]
{匀速滑动时，$I_{1}=0$，$I_{2}=0$}
{匀速滑动时，$I_{1}\neq 0$，$I_{2}\neq 0$}
{加速滑动时，$I_{1}=0$，$I_{2}=0$}
{加速滑动时，$I_{1}\neq 0$，$I_{2}\neq 0$}

%% answer: D
%% memo:
\memoanswer{
电容器在电路中与等效电源并联，两端电压等于 AB 端感应电动势，所以当导体棒匀速滑动时，电容器两端电压不变，所以 $I_{2}=0$，电阻 R 中电流不为零，故 A、B 错误；加速滑动时，电容器两端电压随导体棒速度的增大而增加，电容器持续充电，充电电流不为零，通过电阻的电流也不为零，故 C 错误、D 正确。
}

\item
%% number: 18
%% paperName: 2002年全国理综旧课程第 18 题 6 分
%% typeId: 1
%% type: 单选题
%% chapter: 运动和力
%% point: 牛顿定律中的图象
%% method: 无
%% score: 6
%% degree: 550
%% duplicateId: 0
%% body:
质点所受的力 $F$ 随时间变化的规律如图所示，力的方向始终在一直线上。已知 $t=0$ 时质点的速度为零。在图示 $t_{1}$、$t_{2}$、$t_{3}$ 和 $t_{4}$ 各时刻中，哪一时刻质点的动能最大？\xzanswer{B}

\onepicture[w=5.5cm]{figs/18.png}

\fourchoices[answer=B]
{$t_{1}$}
{$t_{2}$}
{$t_{3}$}
{$t_{4}$}

%% answer: B
%% memo:
\memoanswer{
由力的图像分析可知，在时间 $0\sim t_{1}$ 内，质点向正方向做加速度增大的加速运动；在时间 $t_{1}\sim t_{2}$ 内，质点向正方向做加速度减小的加速运动；在时间 $t_{2}\sim t_{3}$ 内，质点向正方向做加速度增大的减速运动；在时间 $t_{3}\sim t_{4}$ 内，质点向正方向做加速度减小的减速运动，$t_{4}$ 时刻速度为零，因此 $t_{2}$ 时刻质点的速度最大，故 B 正确。
}

\item
%% number: 19
%% paperName: 2002年全国理综旧课程第 19 题 6 分
%% typeId: 1
%% type: 单选题
%% chapter: 光学
%% point: 光的折射
%% method: 无
%% score: 6
%% degree: 550
%% duplicateId: 0
%% body:
为了观察门外情况，有人在门上开一小圆孔，将一块圆柱形玻璃嵌入其中，圆柱体轴线与门面垂直。从圆柱底面中心看出去，可以看到的门外入射光线与轴线间的最大夹角称做视场角。已知该玻璃的折射率为 $n$，圆柱长为 $l$，底面半径为 $r$，则视场角是\xzanswer{B}

\onepicture[w=3.0cm]{\includesvg{figs/19}}

\fourchoices[answer=B]
{$\arcsin\frac{nl}{\sqrt{r^{2}+l^{2}}}$}
{$\arcsin\frac{nr}{\sqrt{r^{2}+l^{2}}}$}
{$\arcsin\frac{r}{n\sqrt{r^{2}+l^{2}}}$}
{$\arcsin\frac{l}{n\sqrt{r^{2}+l^{2}}}$}

%% answer: B
%% memo:
\memoanswer{
观察门外情况时，作出光路图如图所示，

\onepicture[w=3.5cm]{figs/19_an.jpg}

由几何关系得 $\sin\alpha=\frac{r}{\sqrt{r^{2}+l^{2}}}$，由光的折射规律得 $n=\frac{\sin\theta}{\sin\alpha}$，所以视场角 $\theta=\arcsin\frac{nr}{\sqrt{r^{2}+l^{2}}}$，故 B 正确。
}

\item
%% number: 20
%% paperName: 2002年全国理综旧课程第 20 题 6 分
%% typeId: 1
%% type: 单选题
%% chapter: 电路
%% point: 动态分析
%% method: 无
%% score: 6
%% degree: 650
%% duplicateId: 0
%% body:
在如图所示的电路中，$R_{1}$、$R_{2}$、$R_{3}$ 和 $R_{4}$ 皆为定值电阻，$R_{5}$ 为可变电阻，电源的电动势为 $E$，内阻为 $r$。设电流表 A 的读数为 $I$，电压表 V 的读数为 $U$。当 $R_{5}$ 的滑动触点向图中 a 端移动时\xzanswer{D}

\onepicture[w=6.0cm]{figs/20.png}

\fourchoices[answer=D]
{$I$ 变大，$U$ 变小}
{$I$ 变大，$U$ 变大}
{$I$ 变小，$U$ 变大}
{$I$ 变小，$U$ 变小}

%% answer: D
%% memo:
\memoanswer{
当变阻器 $R_{5}$ 滑动触点向 a 端移动时，变阻器接入电路的电阻减小，电路中的总电阻减小，总电流增大，电源内电压增大，路端电压减小，故电压表测量路端电压示数 $U$ 减小，$R_{1}$、$R_{3}$ 两端电压随总电流增大而增大，所以外电路并联部分电压减小，通过电流表所在支路的电流 $I$ 减小，电流 $I$ 示数减小，故 D 正确。
}

\item
%% number: 26
%% paperName: 2002年全国理综旧课程第 26 题 20 分
%% typeId: 6
%% type: 计算题
%% chapter: 运动和力
%% point: 牛顿定律的应用
%% method: 无
%% score: 20
%% degree: 650
%% duplicateId: 0
%% body:
蹦床是运动员在一张绷紧的弹性网上蹦跳、翻滚并做各种空中动作的运动项目。一个质量为 $60\Ukg$ 的运动员，从离水平网面 $3.2\Um$ 高处自由下落，着网后沿竖直方向蹦回到离水平网面 $5.0\Um$ 高处。已知运动员与网接触的时间为 $1.2\Us$。若把在这段时间内网对运动员的作用力当作恒力处理，求此力的大小。（$g=10\Umsq$）

%% answer: 
\jdanswer{
$F=1.5\times10^{3}\UN$
}
%% memo:
\memoanswer{
将运动员看做质量为 $m$ 的质点，从 $h_{1}$ 高处下落，刚接触网时速度的大小为 $v_{1}=\sqrt{2gh_{1}}$（向下）①，

弹跳后到达的高度为 $h_{2}$，刚离网时的速度的大小为

$v_{2}=\sqrt{2gh_{2}}$（向上）②，

速度的改变量为

$\Delta v=v_{1}+v_{2}$（向上）③，

以 $a$ 表示加速度，$\Delta t$ 表示接触时间，则

$\Delta v=a\Delta t$ ④，

接触过程中运动员受到向上的弹力 $F$ 和向下的重力 $mg$，由牛顿第二定律得

$F-mg=ma$ ⑤，

由以上各式解得

$F=mg+m\frac{\sqrt{2gh_{2}}+\sqrt{2gh_{1}}}{\Delta t}$ ⑥，

代入数值得 $F=1.5\times10^{3}\UN$ ⑦。
}

\item
%% number: 27
%% paperName: 2002年全国理综旧课程第 27 题 20 分
%% typeId: 6
%% type: 计算题
%% chapter: 磁场
%% point: 洛伦兹力
%% method: 无
%% score: 20
%% degree: 450
%% duplicateId: 0
%% body:
电视机的显像管中，电子束的偏转是用磁偏转技术实现的。电子束经过电压为 $U$ 的加速电场后，进入一圆形匀强磁场区，如图所示。磁场方向垂直于圆面。磁场区的中心为 O，半径为 $r$。当不加磁场时，电子束将通过 O 点而打到屏幕的中心 M 点。为了让电子束射到屏幕边缘 P，需要加磁场，使电子束偏转一已知角度 $\theta$，此时磁场的磁感应强度 $B$ 应为多少？

\onepicture[w=7.5cm, align=right]{figs/27.png}

%% answer: 
\jdanswer{
$B=\frac{1}{r}\sqrt{\frac{2mU}{e}}\tan\frac{\theta}{2}$
}
%% memo:
\memoanswer{
电子在磁场中沿圆弧 ab 运动，圆心为 C，半径为 R，以 $v$ 表示电子进入磁场时的速度，$m$、$e$ 分别表示电子的质量和电量，在加速电场中运动时，由动能定理得

$eU=\frac{1}{2}mv^{2}$ ①，

由洛伦兹力提供向心力得

$evB=\frac{mv^{2}}{R}$ ②，

由几何关系得 $\tan\frac{\theta}{2}=\frac{r}{R}$ ③，

由以上各式解得

$B=\frac{1}{r}\sqrt{\frac{2mU}{e}}\tan\frac{\theta}{2}$ ④

\onepicture[w=4.0cm]{figs/27_an.jpg}
}

\item
%% number: 29
%% paperName: 2002年全国理综旧课程第 29 题 18 分
%% typeId: 4
%% type: 实验题
%% chapter: 气体性质
%% point: 验证玻意耳定律
%% method: 无
%% score: 18
%% degree: 550
%% duplicateId: 0
%% body:
现用“验证玻意耳定律”的仪器来测量大气压强 $p_{0}$。注射器针筒已被固定在竖直方向上，针筒上所标刻度是注射器的容积，最大刻度 $V_{m}=10\UmL$。注射器活塞已装上钩码框架，如图所示。此外，还有一架托盘天平、若干钩码、一把米尺、一个针孔橡皮帽和少许润滑油。

\begin{enumerate}
\item
下面是实验步骤，试填写所缺的②和⑤。

\begin{steps}
\item
用米尺测出注射器针筒上全部刻度的长度 $L$。
\item
\tkanswer[6cm]{称出活塞和钩码框架的总质量 $M$}
\item
把适量的润滑油抹在注射器的活塞上，将活塞插入针筒中，上下拉动活塞，使活塞与针筒的间隙内均匀地涂上润滑油。
\item
将活塞插到适当的位置。
\item
\tkanswer[8cm]{将注射器针筒上的小孔用橡皮帽堵住}
\item
在钩码框架两侧挂上钩码，记下挂上的钩码质量 $m_{1}$。小达到平衡后，记下注射器中空气柱的体积 $V_{1}$。在这个过程中不要用手接触注射器以保证空气柱温度不变。
\item
增加钩码的个数，使钩码的质量增大为 $m_{2}$，达到平衡后，记下空气柱的体积 $V_{2}$。
\end{steps}
\item
求出计算大气压强 $p_{0}$ 的公式。（用已给的和测得的物理量表示）
\end{enumerate}

\onepicture[w=7.5cm, align=right]{figs/29.png}

%% answer: 
\jdanswer{
\begin{enumerate}
	\item
	②称出活塞和钩码框架的总质量 $M$；⑤将注射器针筒上的小孔用橡皮帽堵住

	\item
	$p_{0}=\frac{Lg}{V_{m}}\left(\frac{m_{2}V_{2}-m_{1}V_{1}}{V_{1}-V_{2}}-M\right)$
\end{enumerate}
}
%% memo:
\memoanswer{
（1）②称出活塞和钩码框架的总质量 $M$；⑤将注射器针筒上的小孔用橡皮帽堵住。

（2）由题意知，活塞的横截面积为

$S=\frac{V_{m}}{L}$ ①，

由力学平衡条件得

$p_{1}=p_{0}+\frac{(M+m_{1})g}{S}$ ②，

$p_{2}=p_{0}+\frac{(M+m_{2})g}{S}$ ③，

由玻意耳定律得

$p_{1}V_{1}=p_{2}V_{2}$ ④，

联立解得大气压强

$p_{0}=\frac{Lg}{V_{m}}\left(\frac{m_{2}V_{2}-m_{1}V_{1}}{V_{1}-V_{2}}-M\right)$ ⑤。
}

\item
%% number: 30
%% paperName: 2002年全国理综旧课程第 30 题 27 分
%% typeId: 6
%% type: 计算题
%% chapter: 电场
%% point: 带电粒子的平衡
%% method: 无
%% score: 27
%% degree: 450
%% duplicateId: 0
%% body:
有三根长度皆为 $l=1.00\Um$ 的不可伸长的绝缘轻线，其中两根的一端固定在天花板上的 O 点，另一端分别拴有质量皆为 $m=1.00\times10^{-2}\Ukg$ 的带电小球 A 和 B，它们的电量分别为 $-q$ 和 $+q$，$q=1.00\times10^{-7}\UC$。A、B 之间用第三根线连接起来。空间中存在大小为 $E=1.00\times10^{6}\UNC$ 的匀强电场，场强方向沿水平方向右，平衡时 A、B 球的位置如图所示。现将 O、B 之间的线烧断，由于有空气阻力，A、B 球最后会达到新的平衡位置。求最后两球的机械能与电势能总和与烧断前相比改变了多少。（不计两带电小球间相互作用的静电力）

\onepicture[w=3.5cm, align=right]{\includesvg{figs/30}}

%% answer: 
\jdanswer{
$\Delta E_{p}=6.8\times10^{-2}\UJ$
}
%% memo:
\memoanswer{
图 1 中虚线表示 A、B 球原来的平衡位置，实线表示烧断后重新达到平衡的位置，其中 $\alpha$、$\beta$ 分别表示细线 OA、AB 与竖直方向的夹角。

\onepicture[h=5.0cm]{figs/30_an1.jpg}

A 球受力如图 2 所示，重力 $mg$，竖直向下；电场力 $qE$，水平向左；细线 OA 对 A 的拉力 $T_{1}$，方向如图所示；细线 AB 对 A 的拉力 $T_{2}$，方向如图，由平衡条件得

\onepicture[h=3.2cm]{figs/30_an2.jpg}
\onepicture[h=3.2cm]{figs/30_an3.jpg}

$T_{1}\sin\alpha+T_{2}\sin\beta=qE$ ①，

$T_{1}\cos\alpha=mg+T_{2}\cos\beta$ ②，

B 球受力如图 3 所示，重力 $mg$，竖直向下，电场力 $qE$，水平向右，细线 AB 对 B 的拉力 $T_{2}$，方向如图所示，由平衡条件得

$T_{2}\sin\beta=qE$ ③，

$T_{2}\cos\beta=mg$ ④，

联立以上各式并代入数据，得

$\alpha=0$ ⑤，

$\beta=45^{\circ}$ ⑥，

由此可知，A、B 球重新达到平衡的位置如图 4 所示，与原来位置相比，A 球的重力势能减少了

\onepicture[h=4.0cm]{figs/30_an4.jpg}

$E_{A}=mgl(1-\sin60^{\circ})$ ⑦，

B 球的重力势能减少了

$E_{B}=mgl(1-\sin60^{\circ}+\cos45^{\circ})$ ⑧，

A 球的电势能增加了

$W_{A}=qEl\cos60^{\circ}$ ⑨，

B 球的电势能减少了

$W_{B}=qEl(\sin45^{\circ}-\sin30^{\circ})$ ⑩，

两种势能总和减少了

$W=W_{B}-W_{A}+E_{A}+E_{B}$ \ccnum{⑪}，

代入数据解得

$\Delta E_{p}=6.8\times10^{-2}\UJ$ \ccnum{⑫}。
}

\end{enumerate}

\end{document}
